\documentclass[11pt]{article}

\usepackage[T1]{fontenc}
\usepackage[utf8]{inputenc}
\usepackage[margin=1in]{geometry}
\usepackage{amsmath,amssymb,mathtools}
\usepackage{enumitem}
\usepackage[hidelinks]{hyperref}

\newcommand{\R}{\mathbb{R}}
\newcommand{\SO}{\mathrm{SO}}
\newcommand{\Exp}{\operatorname{Exp}}
\newcommand{\Log}{\operatorname{Log}}
\newcommand{\IoU}{\operatorname{IoU}}
\newcommand{\clip}{\operatorname{clip}}
\newcommand{\median}{\operatorname{median}}
\newcommand{\normalize}{\operatorname{normalize}}
\newcommand{\SmoothLone}{\operatorname{SmoothL1}}

\title{Mathematics of the GigaPose Adaptive Tracking System}
\author{}
\date{}

\begin{document}
\maketitle
\tableofcontents

\section{System overview}

The tracking system has two distinct parts:

\begin{enumerate}
    \item \texttt{run\_tracking}: deterministic detection association,
    hypothesis generation, CAD rendering, pose scoring, local refinement, and
    temporal tracking.
    \item \texttt{train\_recovery}: training of an optional small neural
    network that proposes pose corrections when a track is uncertain or lost.
\end{enumerate}

The main tracker is not trained. It is a deterministic search and optimization
procedure. Only the optional recovery head is learned. Even when the learned
head is used, its corrected pose is rendered and checked by the deterministic
CAD score before it can be selected.

\section{Pose representation}

Every object pose is represented as a camera-from-object transformation

\begin{equation}
T_{co}=
\begin{bmatrix}
R & t\\
0 & 1
\end{bmatrix},
\qquad
R\in\SO(3),\quad
t=[t_x,t_y,t_z]^T\in\R^3.
\end{equation}

The pose transforms a point in CAD/object coordinates into camera coordinates:

\begin{equation}
X_c=RX_o+t.
\end{equation}

Rotation changes are represented by an axis--angle vector
$\omega\in\R^3$:

\begin{equation}
\Delta R=\Exp([\omega]_\times),
\end{equation}

where $[\omega]_\times$ is the skew-symmetric cross-product matrix. The
exponential map converts a vector in the Lie algebra $\mathfrak{so}(3)$ into a
valid rotation in $\SO(3)$. This prevents an update from producing a
non-orthogonal rotation matrix. The implementation is in
\path{tracking/geometry.py}.

\part{Runtime tracking}

\section{Multi-car association}

Before refining poses, the tracker must decide which detection in the current
frame belongs to each existing track. This is an assignment problem, not a
pose-estimation problem.

\subsection{Basic geometric cost}

For track $i$ and detection $j$, the basic cost is

\begin{equation}
C_{ij}=
w_{\mathrm{IoU}}\bigl(1-\IoU(B_i,B_j)\bigr)
+w_c
\min\left(
\frac{\lVert c_i-c_j\rVert_2/d_{\mathrm{image}}}{d_{\max}},
1
\right),
\end{equation}

where

\begin{equation}
d_{\mathrm{image}}=\sqrt{W^2+H^2}.
\end{equation}

$B_i$ and $B_j$ are bounding boxes, while $c_i$ and $c_j$ are their centers.
The default weights are

\begin{equation}
w_{\mathrm{IoU}}=0.65,
\qquad
w_c=0.35.
\end{equation}

A low cost means that a detection overlaps the previous track box and has a
nearby center.

\subsection{The Hungarian algorithm}

All pairwise costs form a matrix. For two tracks and two detections, for
example,

\begin{equation}
C=
\begin{bmatrix}
C_{0,0} & C_{0,1}\\
C_{1,0} & C_{1,1}
\end{bmatrix}
=
\begin{bmatrix}
0.2 & 0.8\\
0.7 & 0.1
\end{bmatrix}.
\end{equation}

The two possible complete assignments have costs

\begin{align}
(0\rightarrow0,\;1\rightarrow1)&:\quad 0.2+0.1=0.3,\\
(0\rightarrow1,\;1\rightarrow0)&:\quad 0.8+0.7=1.5.
\end{align}

The Hungarian algorithm selects the first assignment. More generally, it
solves

\begin{equation}
\min_{x_{ij}}\sum_{i,j}C_{ij}x_{ij}
\end{equation}

subject to one-to-one assignment constraints

\begin{equation}
x_{ij}\in\{0,1\},
\qquad
\sum_jx_{ij}\le1,
\qquad
\sum_ix_{ij}\le1.
\end{equation}

Thus, two tracks cannot claim the same detection and one track cannot claim
two detections. The algorithm finds the globally best assignment under the
provided cost; it does not create identity information that is absent from
that cost.

\subsection{HSV appearance descriptor}

When identity-aware association is enabled, the tracker computes a color
histogram inside each instance mask. RGB pixels are converted to HSV:

\begin{itemize}
    \item $H$: hue, such as red, blue, green, or yellow;
    \item $S$: saturation, or how strongly colored a pixel is;
    \item $V$: value, or brightness.
\end{itemize}

The current implementation uses only $H$ and $S$, with 24 hue bins and 16
saturation bins. The result is a 384-dimensional normalized histogram

\begin{equation}
h(k)\ge0,
\qquad
\sum_kh(k)=1.
\end{equation}

Hue and saturation are used instead of raw RGB because brightness and color
are more entangled in RGB. Omitting $V$ makes the descriptor less sensitive to
illumination, shadow, haze, and exposure changes. This helps when the same car
moves from sunlight into shadow.

The approach also has limitations. Hue is unreliable for gray, black, and
white cars; identical liveries can have almost identical histograms; the
histogram discards spatial layout; and occlusion can change the observed color
distribution.

\subsection{Bhattacharyya appearance distance}

For two normalized histograms $p$ and $q$, the Bhattacharyya coefficient is

\begin{equation}
BC(p,q)=\sum_k\sqrt{p_kq_k}.
\end{equation}

The corresponding distance used for normalized histograms is

\begin{equation}
d_B(p,q)=\sqrt{1-BC(p,q)}.
\end{equation}

Therefore,

\begin{align}
d_B\approx0 &\quad\Longrightarrow\quad
\text{similar color distributions},\\
d_B\approx1 &\quad\Longrightarrow\quad
\text{different color distributions}.
\end{align}

The stored track appearance is updated with an exponential moving average:

\begin{equation}
h_t=\normalize\left(\mu h_{t-1}+(1-\mu)h_{\mathrm{det}}\right),
\qquad \mu=0.85.
\end{equation}

\subsection{Identity-aware association cost}

When all optional signals are available, the cost numerator is

\begin{equation}
\widetilde C_{ij}=
w_{\mathrm{IoU}}(1-\IoU_{\mathrm{box}})
+w_cd_{\mathrm{center}}
+w_{\mathrm{mask}}(1-\IoU_{\mathrm{mask}})
+w_{\mathrm{app}}d_B
+w_{\mathrm{id}}E_{\mathrm{id}},
\end{equation}

and it is normalized by the sum of weights for the available terms:

\begin{equation}
C_{ij}=\frac{\widetilde C_{ij}}{\sum_{k\in\mathcal A}w_k}.
\end{equation}

Here, $E_{\mathrm{id}}=0$ for matching external IDs and $1$ for different
external IDs. Strict external-ID mode can forbid a conflicting match rather
than merely penalizing it.

Identity-aware association is enabled with

\begin{verbatim}
--identity-aware-association
\end{verbatim}

Stable external instance IDs can additionally be used with

\begin{verbatim}
--use-external-ids
\end{verbatim}

External IDs should not be enabled if the detector independently reassigns
instance numbers in every frame.

If two cars have the same color, $d_B$ is nearly equal for both assignments
and contributes little information. The tracker then relies on box geometry,
mask overlap, motion, and any stable external ID. If two identical cars fully
cross or occlude each other and no persistent identity cue exists, the
assignment is fundamentally ambiguous and an identity swap is possible.

The association implementation is in \path{tracking/association.py}.

\section{Temporal pose prediction}

For consecutive poses, define translation and rotation increments as

\begin{align}
\Delta t_k &= t_k-t_{k-1},\\
\Delta\omega_k &= \Log\left(R_kR_{k-1}^T\right).
\end{align}

The tracker estimates velocity over a short history using the component-wise
median:

\begin{equation}
v_t=\median_k(\Delta t_k),
\qquad
v_R=\median_k(\Delta\omega_k).
\end{equation}

The next temporal proposal is

\begin{align}
\widehat t_{k+1}&=t_k+v_t,\\
\widehat R_{k+1}&=\Exp([v_R]_\times)R_k.
\end{align}

The median is used for robustness. Consider one-dimensional translation
increments

\begin{equation}
[0.10,\;0.11,\;0.09,\;2.00]\;\text{m/frame}.
\end{equation}

The mean is $0.575$ m/frame, while the median is $0.105$ m/frame. A single bad
pose therefore corrupts the mean much more strongly. The same issue occurs
when a temporary front/rear flip creates a rotation increment near
$180^\circ$.

The median is less statistically efficient than the mean when all errors are
small and Gaussian, and it can be less smooth. The behavior is configurable:

\begin{verbatim}
"temporal": {
  "window_size": 5,
  "robust_velocity": true
}
\end{verbatim}

Setting \texttt{robust\_velocity} to \texttt{false} replaces the median with
the mean.

\section{Optical-flow proposal}

Features inside the previous object mask are tracked with pyramidal
Lucas--Kanade optical flow. RANSAC fits a 2D partial affine/similarity model

\begin{equation}
u_t\approx Au_{t-1},
\end{equation}

which contains image translation, in-plane rotation $\theta$, and scale $s$.
The transformed object center supplies the pixel displacement. Image scale is
converted approximately into a depth update:

\begin{equation}
\Delta\log z=-\log s,
\qquad
z_{\mathrm{new}}=z_{\mathrm{old}}e^{-\log s}
=\frac{z_{\mathrm{old}}}{s}.
\end{equation}

Thus, when the object becomes larger in the image ($s>1$), the proposed depth
decreases. The in-plane rotation update is

\begin{equation}
R_{\mathrm{new}}=R_z(\theta)R_{\mathrm{old}}.
\end{equation}

Flow confidence is

\begin{equation}
c_{\mathrm{flow}}=\clip\left(
r_{\mathrm{inlier}}
\min\left(\frac{N_{\mathrm{inlier}}}{30},1\right),
0,1
\right).
\end{equation}

The implementation is in \path{tracking/flow.py}.

\section{Adaptive pose-hypothesis generation}

Depending on the tracking state, the candidate set may combine:

\begin{itemize}
    \item temporal propagation;
    \item optical-flow propagation;
    \item the previous pose;
    \item GigaPose top-$K$ predictions;
    \item optional auxiliary predictions;
    \item $180^\circ$ flipped poses;
    \item image-center, log-depth, translation, and rotation offsets.
\end{itemize}

For a pose with translation $t$, its projected center is

\begin{equation}
u=\pi(Kt).
\end{equation}

Given a pixel offset $\Delta u$ and log-depth offset $\Delta\log z$, define

\begin{equation}
z'=z e^{\Delta\log z}
\end{equation}

and

\begin{equation}
t'=z'
\frac{
K^{-1}
\begin{bmatrix}
u+\Delta u\\1
\end{bmatrix}
}{
\left(
K^{-1}
\begin{bmatrix}
u+\Delta u\\1
\end{bmatrix}
\right)_z
}.
\end{equation}

Equivalent candidates are removed when both

\begin{equation}
\lVert t_i-t_j\rVert_2\le\tau_t
\end{equation}

and

\begin{equation}
\left\lVert\Log(R_iR_j^T)\right\rVert_2\le\tau_R.
\end{equation}

Normal tracks follow a cheap local path. Uncertain and lost tracks receive a
larger candidate set, and periodic global GigaPose hypotheses provide a way to
escape accumulated drift. Candidate generation is implemented in
\path{tracking/hypotheses.py}.

\section{CAD alignment evidence}

For each pose hypothesis, the CAD renderer produces a rendered mask $M_r$ and
rendered depth $D_r$. These are compared with the observed detection mask
$M_o$, observed box $B_o$, and observed depth $D_o$.

\subsection{Silhouette and bounding-box errors}

\begin{align}
E_{\mathrm{sil}}&=1-\frac{|M_r\cap M_o|}{|M_r\cup M_o|},\\
E_{\mathrm{bbox}}&=1-\IoU(B_r,B_o).
\end{align}

\subsection{Symmetric edge error}

Let $\partial M$ denote a mask boundary and let $d(x,S)$ be the distance from
point $x$ to set $S$. The symmetric edge error is

\begin{equation}
E_{\mathrm{edge}}=\clip\left(
\frac{
\frac12\left[
\frac{1}{|\partial M_r|}\sum_{p\in\partial M_r}d(p,\partial M_o)
+\frac{1}{|\partial M_o|}\sum_{q\in\partial M_o}d(q,\partial M_r)
\right]
}{d_{\max}},
0,1
\right).
\end{equation}

This is a normalized, truncated, symmetric Chamfer-style distance.

\subsection{Depth error}

Depth is evaluated only on pixels that are valid in both masks and depth maps:

\begin{equation}
\Omega=M_r\cap M_o\cap\{D_r>0\}\cap\{D_o>0\}.
\end{equation}

The robust log-depth error is

\begin{equation}
E_{\mathrm{depth}}=\clip\left(
\frac{
\median_{p\in\Omega}
\left|\log(D_r(p)+\epsilon)-\log(D_o(p)+\epsilon)\right|
}{\log(1+d_{\mathrm{trunc}})},
0,1
\right).
\end{equation}

If depth is unavailable, the depth term and its weight are removed. Missing
depth is not treated as a zero-error match.

\subsection{Measurement and motion errors}

The term called \texttt{feature\_error} in the implementation currently uses
the incoming GigaPose or auxiliary-prediction score:

\begin{equation}
E_{\mathrm{feature}}=1-\clip(s_{\mathrm{measurement}},0,1).
\end{equation}

It is not a new DINO feature comparison. Relative to a temporal reference
pose, the motion error is

\begin{equation}
E_{\mathrm{motion}}=\frac12\left[
\min\left(\frac{\lVert t-t_{\mathrm{motion}}\rVert_2}{s_t},1\right)
+\min\left(
\frac{\theta(RR_{\mathrm{motion}}^T)}{s_R},1
\right)
\right],
\end{equation}

with default scales

\begin{equation}
s_t=1.5\;\mathrm{m},
\qquad
s_R=35^\circ.
\end{equation}

\subsection{Combined candidate score}

The image-evidence error is

\begin{equation}
E_{\mathrm{evidence}}=
w_fE_{\mathrm{feature}}
+w_sE_{\mathrm{sil}}
+w_eE_{\mathrm{edge}}
+w_dE_{\mathrm{depth}}
+w_bE_{\mathrm{bbox}}.
\end{equation}

The full error is

\begin{equation}
E_{\mathrm{total}}=E_{\mathrm{evidence}}+w_mE_{\mathrm{motion}}.
\end{equation}

Default weights are

\begin{equation}
(w_f,w_s,w_e,w_d,w_b,w_m)
=(0.15,2.0,0.75,0.5,0.75,0.12).
\end{equation}

Candidates are ranked by increasing $E_{\mathrm{total}}$. Visual confidence
deliberately excludes the motion prior:

\begin{equation}
c_{\mathrm{visual}}=\clip\left(
1-\frac{E_{\mathrm{evidence}}}
{w_f+w_s+w_e+w_d+w_b},
0,1
\right),
\end{equation}

where unavailable terms and weights are omitted. Excluding motion prevents a
temporally self-consistent but visually wrong track from becoming permanently
overconfident. Scoring is implemented in \path{tracking/scoring.py}.

\section{Derivative-free local refinement}

The refiner performs coordinate search around each beam candidate. It tests

\begin{equation}
u_x\pm\Delta u,
\qquad
u_y\pm\Delta u,
\qquad
\log z\pm\Delta\log z,
\end{equation}

and, by default, in-plane rotations

\begin{equation}
R_z(\pm\Delta\theta)R.
\end{equation}

Full rotation refinement can additionally search the $x$ and $y$ axes. At
iteration $k$, the search step decays according to

\begin{equation}
\Delta_{k+1}=\gamma\Delta_k,
\qquad
\gamma=0.55\quad\text{by default}.
\end{equation}

Only the lowest-error beam candidates survive. This resembles coordinate
descent with beam search, but it optimizes a CAD-rendering objective without
back-propagating through the renderer. See \path{tracking/refinement.py}.

\section{Confidence, state transitions, and global recovery}

If the two best candidates have $E_1\le E_2$, their separation reliability is

\begin{equation}
r=\clip\left(\frac{E_2-E_1}{m_{\min}},0,1\right).
\end{equation}

Final confidence is

\begin{equation}
c=c_{\mathrm{visual}}(0.75+0.25r).
\end{equation}

An ambiguous result can therefore lose up to 25 percent of its visual
confidence. Default state transitions are

\begin{equation}
\begin{cases}
\text{normal}, & c\ge0.67,\\
\text{uncertain}, & 0.35\le c<0.67,\\
\text{lost}, & c<0.35.
\end{cases}
\end{equation}

Uncertain and lost states perform broader searches. Periodic global
hypotheses prevent unbounded drift. Optional same-frame recovery escalates
from the cheap normal search to uncertain/lost search before the current frame
is finalized. State management is implemented in \path{tracking/tracker.py}.

\part{Learned recovery head}

\section{Recovery-data generation}

Recovery training begins with a ground-truth pose

\begin{equation}
T^*=[R^*,t^*].
\end{equation}

The data generator creates synthetic candidate poses
$\widetilde T=[\widetilde R,\widetilde t]$, including clean poses,
$180^\circ$ flips, moderate and large image-center errors, log-depth errors,
full 3D rotation errors, and partial synthetic occlusions.

Each candidate is rendered and summarized by the 11-dimensional vector

\begin{equation}
x=\left[
E_f,E_{\mathrm{sil}},E_{\mathrm{edge}},E_{\mathrm{depth}},
E_{\mathrm{bbox}},E_{\mathrm{motion}},
\Delta u_x,\Delta u_y,
\log(A_r/A_o),
\log(D_r/D_o),
s_{\mathrm{measurement}}
\right].
\end{equation}

Candidates generated from the same physical object instance share a group ID.
Train/validation splitting is performed by group so perturbations of the same
instance cannot leak across the split. Generation is implemented in
\path{tracking/generate_recovery_dataset.py}.

\section{Recovery targets}

The desired left-applied rotation correction is

\begin{equation}
\omega^*=\Log\left(R^*\widetilde R^T\right),
\end{equation}

which satisfies

\begin{equation}
R^*=\Exp([\omega^*]_\times)\widetilde R.
\end{equation}

The desired camera-frame translation correction is

\begin{equation}
\Delta t^*=t^*-\widetilde t.
\end{equation}

The binary acceptable-pose target is

\begin{equation}
y=\mathbf 1\left[
\lVert t^*-\widetilde t\rVert_2\le0.15\;\mathrm{m}
\;\land\;
E_R\le10^\circ
\right].
\end{equation}

The scalar quality target is

\begin{equation}
q^*=\frac{\lVert t^*-\widetilde t\rVert_2}{0.15\;\mathrm{m}}
+\frac{E_R}{10^\circ}.
\end{equation}

Lower $q^*$ indicates a better candidate.

\section{Recovery network}

Input features are standardized using statistics computed only from the
training groups:

\begin{equation}
\overline x=\frac{x-\mu_x}{\sigma_x}.
\end{equation}

With the default hidden dimension, the multilayer perceptron is

\begin{equation}
11\longrightarrow128\longrightarrow128\longrightarrow8.
\end{equation}

It predicts

\begin{equation}
(\widehat\omega_{\mathrm{raw}},
\Delta\widehat t_{\mathrm{raw}},
\widehat a,
\widehat q),
\end{equation}

comprising three raw rotation values, three raw translation values, a
confidence logit, and a scalar log-quality estimate.

An unconstrained vector is mapped to an open ball of radius $r_{\max}$ by

\begin{equation}
B(v;r_{\max})=
r_{\max}\tanh(\lVert v\rVert_2)
\frac{v}{\max(\lVert v\rVert_2,\epsilon)}.
\end{equation}

Hence

\begin{align}
\widehat\omega&=B(\widehat\omega_{\mathrm{raw}};\omega_{\max}),\\
\Delta\widehat t&=B(\Delta\widehat t_{\mathrm{raw}};t_{\max}).
\end{align}

This prevents the learned head from applying arbitrarily large corrections.
The network is defined in \path{tracking/recovery.py}.

\section{Recovery training objective}

For scalar residual $r$, PyTorch's default Smooth-L1 penalty is

\begin{equation}
\rho(r)=
\begin{cases}
\frac12r^2, & |r|<1,\\
|r|-\frac12, & |r|\ge1.
\end{cases}
\end{equation}

The normalized rotation and translation losses are

\begin{align}
L_R&=\SmoothLone\left(
\frac{\widehat\omega}{\omega_{\max}},
\frac{\omega^*}{\omega_{\max}}
\right),\\
L_t&=\SmoothLone\left(
\frac{\Delta\widehat t}{t_{\max}},
\frac{\Delta t^*}{t_{\max}}
\right).
\end{align}

The binary confidence loss is

\begin{equation}
L_c=-y\log\sigma(\widehat a)
-(1-y)\log\left(1-\sigma(\widehat a)\right).
\end{equation}

The quality loss predicts a compressed log-quality target:

\begin{equation}
L_q=\SmoothLone\left(\widehat q,\log(1+q^*)\right).
\end{equation}

For candidates $i$ and $j$ from the same object instance, if
$q_i^*<q_j^*$, candidate $i$ should have lower predicted quality. The pairwise
ranking penalty is

\begin{equation}
L_{\mathrm{rank}}^{ij}
=\max\left(0,m+\widehat q_i-\widehat q_j\right),
\qquad m=0.1\quad\text{by default}.
\end{equation}

The total objective is

\begin{equation}
\boxed{
L=\lambda_RL_R+\lambda_tL_t+\lambda_cL_c
+\lambda_qL_q+\lambda_{\mathrm{rank}}L_{\mathrm{rank}}
}.
\end{equation}

Default weights are

\begin{equation}
(\lambda_R,\lambda_t,\lambda_c,\lambda_q,\lambda_{\mathrm{rank}})
=(1,1,0.25,0.20,0.20).
\end{equation}

The head is optimized using AdamW with default weight decay $10^{-4}$.
Gradient norm is clipped to 5. Validation uses the same total objective and
early stopping monitors validation $L$. Training and validation loss
components can optionally be logged to Weights \& Biases. See
\path{tracking/train_recovery.py}.

\section{Using the learned correction at runtime}

The learned head runs only when a track is uncertain or lost and a recovery
checkpoint is supplied. For a candidate pose,

\begin{align}
R_{\mathrm{corrected}}
&=\Exp([\widehat\omega]_\times)R_{\mathrm{candidate}},\\
t_{\mathrm{corrected}}
&=t_{\mathrm{candidate}}+\Delta\widehat t.
\end{align}

The predicted quality is converted to a score

\begin{equation}
s_q=e^{-\max(\widehat q,0)},
\end{equation}

and the proposal's measurement score becomes

\begin{equation}
s_{\mathrm{new}}=
\frac{s_{\mathrm{old}}+\sigma(\widehat a)+s_q}{3}.
\end{equation}

The corrected proposal is not accepted blindly. It returns to the CAD refiner
and is compared with all other candidates using silhouette, edge, depth, box,
measurement, and motion evidence. The complete self-correcting loop is

\begin{align*}
\text{temporal/global candidates}
&\rightarrow\text{CAD scoring}
\rightarrow\text{learned correction}\\
&\rightarrow\text{CAD refinement}
\rightarrow\text{verified pose}.
\end{align*}

The approach can improve a bad initial GigaPose pose when candidate generation
or the recovery head reaches the correct alignment basin. It cannot guarantee
recovery when all candidates are far from the true pose or when visual
evidence is fundamentally ambiguous.

\section{Relevant implementation files}

\begin{itemize}
    \item Pose geometry: \path{tracking/geometry.py}
    \item Multi-car association: \path{tracking/association.py}
    \item Optical flow: \path{tracking/flow.py}
    \item Candidate generation: \path{tracking/hypotheses.py}
    \item CAD scoring: \path{tracking/scoring.py}
    \item Local refinement: \path{tracking/refinement.py}
    \item State machine: \path{tracking/tracker.py}
    \item Recovery network: \path{tracking/recovery.py}
    \item Recovery-data generation:
    \path{tracking/generate_recovery_dataset.py}
    \item Recovery training: \path{tracking/train_recovery.py}
    \item Command-line runtime: \path{tracking/run_tracking.py}
\end{itemize}

\end{document}
